%!TEX TS-program = pdftex
% JUAN LUIS VARONA, 14-4-1991.
% RESUMEN DE LA TESIS DOCTORAL PARA LA REVISTA Zubia, DEL I.E.R.
% REFERENCIA A LA REVISTA A\~NADIDA EL 21-8-2004.

\magnification=1200
\pretolerance=500
\tolerance=1000
\baselineskip=18pt % Espacio y medio
\frenchspacing % Para no dejar doble espacio despues de los puntos


\def\'#1{\if#1i{\accent"13\i}\else{\accent"13 #1}\fi}

\centerline{\bf CONVERGENCIA EN $L^p$ CON PESOS DE LA SERIE DE FOURIER} 
\centerline{\bf RESPECTO DE ALGUNOS SISTEMAS ORTOGONALES\rm\footnote{*}{Resumen
de la tesis doctoral; dicha tesis fue dirigida por el Dr. D. Jos\'e J. Guadalupe
y defendida en la Universidad de Cantabria el d\'ia 10-12-1988.
El Instituto de Estudios Riojanos public\'o este resumen
en {\it Zub\'{\i}a\/} {\bf 8} (1990), 217--218.}}

\bigskip


\centerline{\bf Juan Luis Varona}

\centerline{\sl Departamento de Matem\'atica Aplicada. Colegio Universitario de
La Rioja.}


\bigskip


Dada $d\mu$ una medida positiva y $\{\phi_n\}$, $n=0,1,2,\dots$, un sistema
ortogonal completo en $L^2(d\mu)$, para cada $f\in L^2(d\mu)$ se denota $S_nf$ a
las sumas parciales de la serie de Fourier de $f$ respecto del sistema
$\{\phi_n\}$.

Es bien conocido que $S_nf$ converge a $f$ en la norma de $L^2(d\mu)$, pero no
se puede asegurar en general que haya convergencia en $L^p(d\mu)$ para $p$
distinto de 2. Sin embargo, para diversos sistemas ortogonales entre los que se
encuentran los polinomios ortogonales cl\'asicos y los sistemas de Bessel y
Dini s\'i que se conocen condiciones necesarias y suficientes para esta
convergencia.

Un paso m\'as nos lleva a pensar en estudiar la convergencia de $S_nf$ a $f$ en
$L^p(Ud\mu)$ para alguna funci\'on positiva (peso) $U$. M\'as en general a\'un,
dados dos pesos $U$ y $V$, podemos preguntarnos bajo qu\'e condiciones es
cierta la convergencia de $(S_nf)U$ a $fU$ en la norma de $L^p(d\mu)$ para cada
funci\'on $f$ tal que $fV\in L^p(Ud\mu)$, $1<p<\infty$. Cuando $U\le CV$ esto
se traduce en la acotaci\'on uniforme de $S_n$ como operador entre
$L^p(V^pd\mu)$ y $L^p(U^pd\mu)$.

Para realizar este estudio es fundamental disponer de buenas estimaciones de
las funciones ortogonales y de una adecuada descomposici\'on en sumandos del
n\'ucleo de la serie de Fourier. De este modo el problema se reduce,
fundamentalmente, a conocer la acotaci\'on de la transformada de Hilbert entre
diversos espacios de tipo $L^p$ con pesos, para lo cual se emplea la teor\'ia
$A_p$ de pesos.

En el trabajo se analizan todos estos problemas para diversos sistemas
ortogonales que incluyen $a)$ polinomios de Jacobi generalizados, $b)$ sistemas
provinientes de pesos de pesos de Jacobi a los que se han a\~nadido deltas de
Dirac, $c)$ sistemas de Bessel y Dini, $d)$ polinomios de Hermite
generalizados. Explicamos a continuaci\'on brevemente lo que se obtiene en
relaci\'on con cada uno de los apartados anteriores.

Para pesos en el intervalo $(-1,1)$ en los que los polinomios ortonormales
verifican determinadas acotaciones, como es el caso de los polinomios de Jacobi
generalizados, se encuentran condiciones necesarias y suficientes para la
convergencia en media deseada. Estas condiciones se expresan en t\'erminos de
$p$ y de condiciones $A_p$ y de integrabilidad para los diversos pesos que
intervienen. Hay que hacer notar que cuando el peso de Jacobi es
$(1-x)^\alpha(1+x)^\beta$ con $\alpha,\beta\ge-1/2$ todas las condiciones se
expresan de forma m\'as sencilla, y las demostraciones son m\'as f\'aciles de
hacer. Otro sistema ortogonal en $(-1,1)$ que se analiza es el de Pollaczek,
para el cual se muestra que no puede haber convergencia en media para $p\not=2$.

Para estudiar la convergencia de sistemas provinientes de pesos de Jacobi con
deltas de Dirac es necesario previamente calcular estimaciones de los
polinomios ortonormales, pues no eran conocidas. Es de destacar que, hasta
donde el autor conoce, esta es la primera vez que se aborda la convergencia de
series de Fourier de polinomios ortogonales respecto a una medida no
absolutamente continua.

En los sistemas de Bessel y Dini --que ya no est\'an constituidos por
polinomios-- se encuentran condiciones necesarias y condiciones suficientes
para la convergencia en media que, como ocurr\'ia con polinomios de Jacobi, se
expresan en funci\'on de $p$ y de condiciones $A_p$ y de integrabilidad sobre los
pesos implicados. De nuevo en este caso cuando el par\'ametro de la funci\'on de
Bessel es $\alpha\ge-1/2$ las condiciones tienen una expresi\'on m\'as sencilla y
las demostraciones se simplifican.

Por \'ultimo, el m\'etodo empleado para analizar la convergencia de series de
Fourier de polinomios de Hermite generalizados es radicalmente diferente del
utilizado para los otros sistemas. Consiste simplemente en usar resultados
conocidos sobre series de polinomios de Laguerre y adaptarlos para obtener
conclusiones sobre series de polinomios de Hermite generalizados. El m\'etodo
resulta sencillo de aplicar y los resultados que se consiguen permiten incluso
extender lo ya conocido para el sistema original de Hermite.

\end
